% Part: sets-functions-relations
% Chapter: infinite
% Section: dedekind-induction

\documentclass[../../../include/open-logic-section]{subfiles}

\begin{document}

\olfileid[ar]{sfr}{infinite}{induction}
\olsection[الاستقراء الحسابي]{جبرات ديديكند والاستقراء الحسابي}

النقطة الحاسمة الآن هي أن جبر ديديكند---بل \emph{كل} جبر
ديديكند---يصلح بديلًا من الأعداد الطبيعية. ويرجع ذلك إلى النتيجة
البسيطة الآتية:

\begin{thm}[الاستقراء الحسابي]\ollabel{thm:dedinfiniteinduction}
لتكن $N, s, o$ مكوّنات جبر ديديكند. عندئذ، لكل مجموعة $X$:  
	\begin{center}
		إذا كان $o \in X$ وكان $(\forall {n} \in N \cap X){s}({n}) \in X$، {فإن} $N \subseteq X$.
	\end{center}
\end{thm}

\begin{proof}
بحسب تعريف جبر ديديكند، $N = \closureofunder{s}{o}$.
والآن، إذا كان كل من ${o} \in X$ و$(\forall {n} \in N)(n \in X \lif
{s}({n}) \in X)$، فإن $N = \closureofunder{s}{o} \subseteq X$.
\end{proof}

وبما أن الاستقراء من الخواص المميزة للأعداد الطبيعية، فالمقصود هو هذا.
إذا أعطيت أي مجموعة لانهائية بحسب ديديكند، أمكننا تكوين جبر ديديكند
واستعمال ذلك الجبر بديلًا من الأعداد الطبيعية.

لا شك أن \olref{thm:dedinfiniteinduction} تصوغ الاستقراء بعبارات
\emph{نظرية-مجموعاتية}. لكن يسهل علينا وضع المبدأ في عبارات قد تكون
أكثر ألفة:

\begin{cor}\ollabel{natinductionschema}
لتكن $N, s, o$ مكوّنات جبر ديديكند. عندئذ، لكل صيغة
$\phi(x)$ قد تكون لها معلمات:
\begin{center}
	إذا كانت $\phi(o)$ وكان $(\forall {n} \in N)(\phi(n)\lif
	\phi({s}({n})))$، {فإن} $(\forall n \in N)\phi(n)$
\end{center}
\end{cor}

\begin{proof}
لتكن $X = \Setabs{n \in N}{\phi(n)}$، واستعمل الآن
\olref{thm:dedinfiniteinduction}
\end{proof}

تكلمنا في هذه النتيجة على صيغة «لها معلمات». ومعنى ذلك، على وجه
التقريب، أنه يمكننا، لأي أشياء $c_1, \ldots, c_k$، أن نعمل بالصيغة
$\phi(x, c_1, \ldots, c_k)$. وبتعبير أدق، يمكننا تقرير النتيجة من غير
ذكر «المعلمات» كما يأتي. لكل صيغة $\phi(x, v_1, \ldots, v_k)$ تعرض
جميع متغيراتها الحرة، لدينا:
	\begin{align*}
			\forall v_1 \ldots \forall v_k((&\phi(o, v_1,\ldots, v_k) \land {}\\
			&	(\forall x \in N)(\phi(x,v_1, \ldots, v_k) \lif \phi(s(x), v_1,\ldots, v_k))) \lif {}\\
			&\hspace{3em} (\forall x \in N)\phi(x, v_1,\ldots, v_k))
	\end{align*}
ومن الواضح أن الكلام على «وجود معلمات» قد يجعل العبارات أيسر قراءة.
(وسنستعمل هذا الأسلوب كثيرًا في \olref[sth][][]{part}.)

وبالعودة إلى جبرات ديديكند: يمكننا، إذا أعطينا أي جبر ديديكند، أن نعرّف
أيضًا الدوال الحسابية المعتادة، وهي الجمع والضرب والأس. لكن ذلك غير
بسيط، وينطوي على تقنية \emph{التعريف العودي}. وهذه تقنية سندخلها
ونسوغها في موضع متأخر كثيرًا وفي سياق أعم بكثير. (وقد يرغب المتحمسون
في العودة إلى هذا بعد \olref[sth][ord-arithmetic][]{chap}، أو لعلهم
يقرؤون معالجة بديلة، مثل \citealt[صص.~95--8]{Potter2004}.) ولكن حيث
تكون $N, s, o$ مكوّنات جبر ديديكند، سيكون في مقدورنا في نهاية المطاف أن نشترط
ما يأتي:
\begin{align*}
	{a} + {o} &= {a} & & & {a} \times {o} &= {o} & & & {a}^{o} &= s(o)\\
{a} + {s}({b}) &= {s}({a}+{b}) &&& {a} \times {s}({b}) &= ({a}\times {b}) + {a}  & & & {a}^{{s}({b})} &= {a}^{b} \times {a}
\end{align*}
وأن نبين أن هذه الدوال تسلك كما نأمل.

\end{document}
